\textit{Przypominam: Zajęcia grupy 2. 17.12.2025 zaczynają się wyjątkowo o 12:00}

\section{Zadanie 1}
Udowodnij, że rozkład wykładniczy ma uogólnioną własność
bez pamięci. To znaczy, jeśli \(X\) jest zmienną o rozkładzie wykładniczym,
\(Y\) i \(Z\) są zmiennymi o rozkładach ciągłych przyjmującymi
wartości nieujemne, i \(X, Y, Z\) są niezależne to mamy
\[ P(X>Y+Z \mid X>Y) = P(X>Z) \]

\subsection*{Rozwiązanie}

Załóżmy, że paramatrem rozkładu \(X\) jest \(\lambda\)
(\( X \sim \text{Exp}(\lambda)\))

\[
    P(X>Z) = \int_0^{\infty} f_Z(z) \cdot (1-F_X(z))~dz
    = \int_0^{\infty} f_Z(z) e^{-\lambda z}~dz
\]
--- dla każdego nieujemnego \(z\) rozważamy zdarzenie, że \(X\) jest od niego
większy.

Rozpisując prawdopodobieństwo warunkowe:

\[
    P(X>Y+Z \mid X>Y) = \frac{P(X>Y+Z~\land~X>Y)}{P(X>Y)}
    = \frac{P(X>Y+Z)}{P(X>Y)}
\]
\[
    P(X > Y+Z)
    = \int_0^{\infty} \int_0^{\infty} f_Y(y) f_Z(z) 
    \cdot (1-F_X(y+z))~dy~dz
    = \int_0^{\infty} \int_0^{\infty} f_Y(y) f_Z(z) 
    e^{-\lambda(y+z)}~dy~dz
\]
\[
    = \int_0^{\infty} f_Y(y) e^{-\lambda y}~dy
    \int_0^{\infty} f_Z(z) e^{-\lambda z}~dz
    = P(X>Y) \cdot P(X>Z)
    \quad \square
\]
